% =====================================================================
% Slide 2: Motivation -- engineering significance of (E, nu)
% =====================================================================
\begin{frame}{Motivation: Elastic Constants as Primary Design Criteria}
\begin{columns}[T]
\column{0.42\textwidth}
\textbf{Two elastic constants govern the structural response of
every load-bearing component:}

\smallskip
\begin{itemize}\small
  \item $E = \sigma/\varepsilon$ \\
        \scriptsize Young's modulus --- uniaxial stiffness (GPa).
  \item $\nu = -\varepsilon_{\text{lat}}/\varepsilon_{\text{axial}}$ \\
        \scriptsize Poisson's ratio --- transverse contraction.
\end{itemize}

\smallskip
For cubic crystals these derive analytically from the two
independent stiffness constants $C_{11}$ and $C_{12}$ ---
the quantities targeted from first principles in this work.

\column{0.57\textwidth}
\begin{block}{Representative structural applications}
\scriptsize
\setlength{\tabcolsep}{4pt}
\renewcommand{\arraystretch}{1.05}
\begin{tabular}{p{2.2cm}lrr}
\toprule
Application & Material & $E$ (GPa) & $\nu$ \\
\midrule
Aerospace fuselage    & Al\,7075         & 71--72 & 0.33 \\
Turbine disk          & Inconel 718      & 200    & 0.29 \\
Orthopaedic implant   & Ti\,6Al--4V      & 110    & 0.34 \\
Automotive structure  & Mg AZ91          & 45     & 0.35 \\
Structural steel      & Mild steel       & 200    & 0.30 \\
MEMS resonator        & Si (single xtal) & 130    & 0.28 \\
Rolling-element bearing & 52100 steel    & 210    & 0.30 \\
\bottomrule
\end{tabular}
\end{block}
\end{columns}
\end{frame}

% =====================================================================
% Slide 3: (E, nu) in mechanical design
% =====================================================================
\begin{frame}{Role of $(E, \nu)$ in Mechanical Design}
\begin{columns}[T]
\column{0.50\textwidth}
\textbf{Quantities governed by $E$:}
\begin{itemize}\scriptsize
  \item Beam deflection $\delta = FL^{3}/(3EI)$
  \item Euler buckling $P_{\mathrm{cr}} = \pi^{2}EI/(KL)^{2}$
  \item Natural frequency $\omega_{n} \propto \sqrt{E/\rho}$
  \item Hertzian contact pressure $\propto E^{*}$
\end{itemize}

\smallskip
\textbf{Quantities governed by $\nu$:}
\begin{itemize}\scriptsize
  \item Bulk modulus $K = E/[3(1-2\nu)]$
  \item Plane-stress vs.\ plane-strain regimes
  \item Stress concentration at geometric discontinuities
  \item Fracture toughness $\propto (1-\nu^{2})$
\end{itemize}

\column{0.49\textwidth}
\begin{block}{Representative ranges}
\scriptsize
\setlength{\tabcolsep}{4pt}
\begin{tabular}{lrr}
\toprule
Material           & $E$ (GPa) & $\nu$ \\
\midrule
Polymers           & 0.5 -- 3  & 0.40 \\
Mg alloys          & 40 -- 45  & 0.35 \\
\textbf{Al\,7075}  & \textbf{70 -- 75} & \textbf{0.33} \\
Ti 6Al--4V         & 110       & 0.34 \\
Steel              & 200       & 0.30 \\
Rubber             & $\sim$0.01 & $\to 0.5$ \\
\bottomrule
\end{tabular}
\end{block}

\smallskip
\begin{block}{Mechanical stability criteria}
\scriptsize
Isotropic continuum: $-1 < \nu < 0.5$.\\
Cubic crystal (Born--Cauchy): $C_{11} > C_{12}$.
\end{block}
\end{columns}
\end{frame}
